\chapter{Specification and Design}%
\label{chp:design}

\section{System Overview}%
\label{sec:design_overview}

The system is a modular computer vision pipeline that accepts broadcast football video clips as input and produces, for each annotated offside event, a binary verdict of offside or onside together with the separation in metres between the attacker and the second-to-last defender. The pipeline is decomposed into six sequential stages, shown in Figure~\ref{fig:pipeline_overview}. Each stage has well-defined inputs and outputs, which allows individual stages to be developed, tested, and replaced independently.

\begin{figure}[htbp]
\centering
\begin{tikzpicture}[
  node distance = 0.55cm,
  auto/.style  = {rectangle, draw, rounded corners=4pt,
                  fill=blue!8, text width=7.2cm,
                  align=center, minimum height=0.9cm, font=\small},
  man/.style   = {rectangle, draw, rounded corners=4pt,
                  fill=orange!12, text width=7.2cm,
                  align=center, minimum height=0.9cm, font=\small},
  arr/.style   = {-{Stealth[length=5pt]}, thick}
]
\node[auto] (acq)  {Stage 1 — Data Acquisition\\[-1pt]
                    {\footnotesize SoccerNet broadcast clips}};
\node[man,  below=of acq]  (kf)   {Stage 2 — Keyframe Selection\\[-1pt]
                    {\footnotesize Manual annotation interface}};
\node[auto, below=of kf]   (det)  {Stage 3 — Player Detection \& Team Identification\\[-1pt]
                    {\footnotesize YOLO11 segmentation + K-Means clustering}};
\node[man,  below=of det]  (lm)   {Stage 4 — Landmark Picking\\[-1pt]
                    {\footnotesize Manual annotation interface}};
\node[auto, below=of lm]   (hom)  {Stage 5 — Homography Estimation\\[-1pt]
                    {\footnotesize RANSAC + least-squares refit}};
\node[man,  below=of hom]  (off)  {Stage 6 — Offside Decision\\[-1pt]
                    {\footnotesize Geometric verdict interface}};
\draw[arr] (acq) -- (kf);
\draw[arr] (kf)  -- (det);
\draw[arr] (det) -- (lm);
\draw[arr] (lm)  -- (hom);
\draw[arr] (hom) -- (off);
\end{tikzpicture}
\caption{Pipeline overview. Blue stages are fully automated; orange stages require operator input.}
\label{fig:pipeline_overview}
\end{figure}

The three automated stages (1, 3, and 5) run without human interaction and produce deterministic output given fixed inputs. The three manual stages (2, 4, and 6) are implemented as lightweight browser-based interfaces served by Python's built-in HTTP server, allowing the operator to interact with the pipeline from any browser without installing additional software. The deliberate inclusion of manual stages reflects the scope constraints stated in Section~\ref{sec:intro_scope}: automating keyframe selection and landmark detection are non-trivial problems in their own right, and removing them from scope allows the evaluation to focus on the accuracy of the geometric and classification components.

\section{Pipeline Stages}%
\label{sec:design_stages}

\subsection{Stage 1 — Data Acquisition}%
\label{sec:design_acq}

The input data is drawn from the SoccerNet dataset~\cite{giancola2018soccernet,deliegev2soccernet}, which provides broadcast video of professional football matches together with temporal annotations for notable events including offside decisions. Match videos are downloaded using the SoccerNet Python package, and a clip-extraction script trims each annotated offside event to a short window centred on the event timestamp. Each clip is stored as an MP4 file named with the match identifier, half, and event index, producing a collection of short broadcast clips ready for annotation.

\subsection{Stage 2 — Keyframe Selection}%
\label{sec:design_keyframe}

For each clip, the user must identify the exact frame at which the ball leaves the passer's foot, since this is the instant at which the offside condition is evaluated according to Law~11. A browser-based interface displays each extracted clip frame by frame, providing the user with buttons for single-frame stepping, ten-frame and twenty-five-frame jumps, along with a slider at the top for easier browsing. Pressing the mark button appends the current frame index to a JSON annotation file. The interface also supports skipping clips that are unsuitable, this was done because there was a number of clips in which the player who ends up receiving the ball would not be visible at the same time that the ball was played and therefore these inappropriate clips could not be evaulated under my system. The output of this stage is a mapping from each clip filename to the annotated frame index.

\subsection{Stage 3 — Player Detection and Team Identification}%
\label{sec:design_detection}

The annotated frame index obtained from the previous step is now used to export each keyframe as a JPEG image. These images are then passed to the YOLO11 object detector to detect all visible players. The segmentation variant of YOLO11 is used in preference to the plain detection variant, because it produces per-pixel masks for each detected player, which allows kit colour sampling to be restricted to body pixels rather than the rectangular bounding box region. Detection is performed in two passes to maximise recall across the full depth of the pitch. A primary pass at a confidence threshold of 0.3 collects all persons above an adaptive minimum height. A secondary pass at confidence 0.05 is then run; only detections whose bounding box centre falls in the upper portion of the frame --- the far end of the pitch, where perspective compression makes distant players appear small --- are accepted from this pass. Duplicate detections between passes are suppressed by an intersection-over-union overlap check.

For each detection, a torso crop is extracted by restricting the bounding box to the top two thirds of the player's height, done intentionally to reduce contribution of players shoes and socks which often vary in colour and could significantly effect the calculation of the clustering algorithm. Apart from that, the bounding box was also trimmed horizontally to possibly eliminate the arms and sleeves. Grass-coloured pixels are identified by computing the mean grass colour from the image in HSV colour space and excluded before colour analysis. The remaining pixels are converted to CIE Lab colour space and summarised by a six-dimensional feature vector comprising the median L*, a*, and b* values, together with three additional statistics capturing the proportion of dark pixels, the proportion of bright pixels, and a saturation proxy. This richer representation improves separation between dark and light kits that are close in hue.

K-Means clustering with $k = 2$ is then applied to the set of player feature vectors to partition players into two teams. Detections belonging to neither team (referees, partially visible players) may be reassigned or excluded at this stage. The output is a JSON file recording, for each keyframe, the set of player bounding boxes and their assigned team labels.

\subsection{Stage 4 — Landmark Picking}%
\label{sec:design_landmarks}

To estimate the homography between the image plane and the real-world pitch, a set of point correspondences must be established. A second browser-based interface displays each keyframe and allows the operator to click on visible pitch landmarks such as line intersections, corner flags, and penalty-spot markings. For each click, the operator also enters the known real-world coordinates of that landmark in the FIFA standard pitch coordinate system, in which the origin is at the pitch centre and coordinates are expressed in metres. Each click is refined to subpixel accuracy before being stored. A minimum of four correspondences are required per keyframe, since each correspondence contributes two linear constraints on $\mathbf{H}$ and the homography has eight degrees of freedom, making four the mathematical minimum for a unique solution (as derived in Section~\ref{sec:bg_geometry}); in practice, between five and ten are typically identified to over-constrain the system and improve robustness.

\subsection{Stage 5 — Homography Estimation}%
\label{sec:design_homography}

The landmark correspondences are passed as input to a homography estimation script, which uses RANSAC to establish a reliable $3 \times 3$ homography matrix $\mathbf{H}$ mapping image coordinates to real-world pitch coordinates. A reprojection threshold of 2.0\,m is used. The threshold must be calibrated to the precision achievable with manual landmark clicks at broadcast-camera scale: one pixel near the far touchline corresponds to approximately 0.06\,m in world space, so a threshold of 0.5\,m would permit only $\pm$8\,pixels of click error, which is stricter than the $\pm$10--15\,pixel localisation precision of a human operator. A threshold of 2.0\,m correctly accepts well-placed clicks at all depths of the visible pitch while still rejecting gross errors --- a misidentified landmark is typically displaced by tens of pixels, corresponding to several metres of world-space error. The final homography is then established using an ordinary least-squares refit on the RANSAC inlier set, which yields a tighter solution than the RANSAC homography alone. The quality of each estimated homography is then reported as the symmetric pixel reprojection error, calculated by mapping the original world coordinates back through $\mathbf{H}^{-1}$ and measuring the Euclidean distance to the clicked point. A validation overlay image is also produced, projecting the FIFA pitch line model into the keyframe so that alignment errors are immediately visible.

\subsection{Stage 6 — Offside Decision}%
\label{sec:design_offside}

The final stage applies the offside rule geometrically. A browser extension provides a user interface which displays each keyframe along with the player detections and their team identifications. It asks the user to choose the attacking team, the direction of attack and the attacker who ends up receiving the ball. The system then automatically identifies the second-to-last defender by projecting the leading edge of each defending player's bounding box through $\mathbf{H}$ into world coordinates and ranks the defenders by their world-space position relative to the direction of attack, the most advanced outfield defender is then selected.

With both offensive and defensive players identified, the system projects their foot positions (the lower midpoint of each bounding box) into world coordinates via $\mathbf{H}$, extracts the component along the direction of attack, and compares the two values. If the attacker's world position is further forward than the defender's, the verdict is offside; otherwise it is onside. The interface also provides the user with the margin between the two players in metres and is reported together with the verdict. An offside line is also back-projected into the image for visual verification. The verdict and all associated data are written to a JSON output file for subsequent evaluation and analysis.

\section{Design Decisions and Alternatives Considered}%
\label{sec:design_decisions}

\subsection{Choice of Object Detector}%
\label{sec:design_det_choice}

YOLO11 was selected as the player detection model because it offers a balance between being accurate and being fast, along with the availability of pretrained weights that cover the \texttt{person} class. Two-stage detectors such as Faster R-CNN~\cite{ren2015faster} were considered but rejected on both speed and accuracy grounds. Benchmarked on the same hardware used in this project (NVIDIA RTX~3060), Faster R-CNN with a ResNet-50-FPN backbone requires approximately 487\,ms per image, compared with 45\,ms for a single YOLO11x-seg pass — roughly a 10$\times$ slowdown. Critically, this speed penalty does not come with an accuracy advantage: Faster R-CNN R-50-FPN achieves approximately 37.4\% box AP on COCO~\cite{ren2015faster}, while YOLO11x-seg achieves 54.7\%~\cite{ultralytics2024yolo11}, making YOLO11x-seg both faster and more accurate on this benchmark. Although inference time is not a hard constraint for a static single-keyframe pipeline, the combination of lower accuracy and 10$\times$ higher latency gives no justification for choosing Faster R-CNN. Of the  one-stage detectors, YOLO11 is a state of the art for the YOLO family~\cite{ultralytics2024yolo11} and is more accurate than its predecessors on the COCO benchmark. The largest segmentation variant, YOLO11x-seg, was selected over smaller variants. YOLO11x achieves the highest reported box mAP of 54.7\% on the COCO benchmark among all YOLO11 variants~\cite{ultralytics2024yolo11}. To confirm this advantage on the actual dataset, YOLO11n-seg, YOLO11m-seg, and YOLO11x-seg were each run at a confidence threshold of 0.3 across all 164 annotated keyframes. YOLO11x-seg produced a mean of 15.3 detections per frame compared with 14.5 for both YOLO11m-seg and YOLO11n-seg, yielding on average 0.8 additional detections per frame and outperforming YOLO11m-seg in the majority of frames. The additional detections are concentrated in distant or partially occluded players at the far end of the pitch, precisely the cases where a smaller model loses recall. This empirical advantage was also observed qualitatively during development, where YOLO11m-seg frequently missed players in the far half of the pitch (see Section~\ref{sec:impl_challenges}). The per-pixel segmentation masks produced by the segmentation variant additionally allow kit colour sampling to be restricted to pixels belonging to the player's body, reducing contamination from background grass and improving team identification accuracy.

\subsection{Team Identification Method}%
\label{sec:design_teamid_choice}

Kit colour clustering in CIE Lab colour space was chosen as the team identification method in preference to several alternatives. A deep feature-based approach, in which a pretrained \gls{cnn} is used to extract appearance embeddings for each player, would require either fine-tuning on labelled football data or a zero-shot embedding model, both of which add considerable complexity and computational cost. Benchmarked on the project hardware, passing a batch of 19 player crops through a ResNet-50 feature extractor requires approximately 37\,ms, compared with 1.6\,ms for K-Means clustering on the same 19 players — a 23$\times$ overhead with no labelled football data to demonstrate a compensating accuracy gain. Colour-based clustering is sufficient because the Laws of the Game require teams to wear distinct kits. CIE Lab was chosen over RGB or HSV because its Euclidean distance is perceptually uniform, producing more stable K-Means boundaries across varying lighting conditions. The six-dimensional feature vector, which adds brightness and saturation statistics to the three Lab channel medians, was found to improve discrimination between kits that are spectrally similar but differ in lightness.

\subsection{Manual Keyframe Selection}%
\label{sec:design_kf_choice}

Keyframe selection is performed manually. Automated ball detection is unreliable in broadcast footage where the ball is frequently occluded or blurred, and an incorrect keyframe directly corrupts the verdict. The operator time cost of 10--30 seconds per event is a justified trade-off for guaranteed correctness.

\subsection{Manual Landmark Picking}%
\label{sec:design_landmark_choice}

Pitch landmark identification is similarly performed manually. Automated approaches such as Hough-based line detection or semantic segmentation are sensitive to player occlusion and variable broadcast angles. Manual picking allows the operator to select the most reliable landmarks per frame, producing higher-quality correspondences than a fixed automated method.

\subsection{Foot-Level Projection and World-Space Defender Ranking}%
\label{sec:design_projection_choice}

The pitch position for each player is the projection of the leading-edge foot point of their bounding box: the bottom corner on the side facing the direction of attack. This follows the geometric interpretation of Law~11, which checks the body part closest to the goal line, and avoids the overestimation that would result from using the centroid.

The most advanced defender is identified by ranking players in world space after projecting through $\mathbf{H}$, avoiding the perspective distortion that makes image-space ranking unreliable on oblique cameras. This adds no computational cost since the homography has already been computed.
